\documentclass[11pt,a4paper]{article}
\usepackage[utf8]{inputenc}
\usepackage[T1]{fontenc}
\usepackage{lmodern,amsmath,amssymb,textcomp}
\usepackage[margin=24mm]{geometry}
\usepackage{longtable,booktabs,array,enumitem}
\usepackage[hidelinks]{hyperref}
\setlength{\parindent}{0pt}
\setlength{\parskip}{6pt}
\emergencystretch=3em
\hypersetup{pdfauthor={Woohyuk Kang, Alejandro Zarzuelo Urdiales}}
\begin{document}
{\LARGE\bfseries Star edge colourings of generalized Petersen graphs and certified reductions towards the six-colour conjecture\par}\medskip
{\small Woohyuk Kang\ensuremath{^1}, Alejandro Zarzuelo Urdiales\ensuremath{^2}\par}\medskip
{\small\textbf{Affiliations:} \href{https://www.kyunghee.edu/eng/user/main/view.do}{\ensuremath{^1} Kyung Hee University, Computer Science and Engineering}; \href{https://rsihouse.ai/openmath}{\ensuremath{^2} Open Math}.\par}

{\small\href{https://github.com/alejandrozu/openmath-2026-judging}{Code and formal proofs}\par}

\subsection{Abstract}
The latest construction gives five-star edge colourings of GP(n,k) for 1\ensuremath{\leq}k\ensuremath{\leq}15 and n\ensuremath{\geq}2k+1, except GP(3,1), alongside other infinite families, small cubic multigraph classifications and conditional reduction theorems. The general six-colour conjecture remains open.

\ensuremath{\chi}\ensuremath{^{\prime}}\_s(GP(n,k))\ensuremath{\leq}5 for 1\ensuremath{\leq}k\ensuremath{\leq}15, n\ensuremath{\geq}2k+1, (n,k)\ensuremath{\neq}(3,1); additionally selected infinite-family colourings, finite six-colour cases and explicitly conditional reduction theorems.

\subsection{Definitions and principal unconditional scope}
A star edge colouring is a proper edge colouring with no two-coloured path or cycle of four edges. Parallel edges in a loopless multigraph are distinct and adjacent. The Dvořák-Mohar-Šámal conjecture asks whether every subcubic multigraph has star chromatic index at most six; the located general published bound is seven.

The strongest current generalized-Petersen endpoint states that GP(n,k) has a five-star edge colouring for 1\ensuremath{\leq}k\ensuremath{\leq}15 and n\ensuremath{\geq}2k+1, except(n,k)=(3,1). Here GP(n,k) has an outer n-cycle, n spokes and inner edges of step k. The condition n>2k makes the construction an ordinary simple cubic graph. This theorem includes parameter regions not covered by the earlier even-order or gcd-based results.

Additional explicit endpoints give colourings of flower and Goldberg snarks, Möbius and inflation families, and a matching-colour-class six-colour construction for GP(n,2). Known parameter subfamilies remain background to these constructions.

\subsection{Periodic blocks, seams and local-window lifting}
The GP(n,k) construction colours repeated blocks and inserts a finite seam determined by the residue of n modulo a selected period. Each k has an explicit ordinary graph colouring recipe, small exceptional tables and representative windows. For example, the k=4 development repeats a period-five block, handles 9\ensuremath{\leq}n<19 directly and compares later windows with representatives 25+(n mod 5).

A forbidden four-edge path or four-cycle touches only a bounded range of blocks. The BlockStar.star\_of\_windows theorem embeds such a candidate into a window of length 3D+1, where D bounds the relevant edge jump. A finite representative-window check therefore rules out every forbidden pattern in all large n. This is an infinite proof supported by finite local certificates, rather than an inference that a finite experimental table probably continues.

The independent ordinary audit regenerated 1,366 GP instances across k=1,...,15, including all small orders through 100 and large examples near 1000, and checked proper colouring and every relevant four-edge walk. Fifty-one flower-snark examples also passed. All GP JSON tables matched the actual Lean literals. These 1,417 tests support semantic fidelity and the finite data, while the formal local-window theorem supplies the intended unbounded conclusion.

\subsection{Small multigraphs and the reduction programme}
The unconditional finite endpoints include six-star edge colourability of all bridgeless cubic loopless multigraphs with at most 14 vertices, and of all subcubic multigraphs with at most seven vertices. They are not interchangeable: lifting a cubic census result to arbitrary subcubic graphs can increase the number of vertices, so the 14-vertex threshold cannot simply be copied.

The reduction library organizes a general loopless-subcubic theorem into cubic and leaf statements. Its basic equivalence says the full conjecture follows from six-colourability of bridgeless cubic graphs together with the six-colourability of their appropriate leaf attachments. Components, bridges and a degree-two core are handled by restriction and gluing lemmas. Two-edge-cut decomposition and digon insertion then concentrate the unresolved cases on more structured cubic hosts.

The library also introduces perfect-matchings-as-colour-classes, far-exchange sets and pole-dominance data. An MC colouring has one prescribed perfect matching as its sixth colour class. EX1-goodness asks for compatible choices preserving an edge's membership status; pole dominance encodes the boundary colour compatibility needed to glue cubic pieces. Theorems prove implications from precisely stated remaining hypotheses. They do not prove those hypotheses globally.

\subsection{Prior work and the actual new regions}
The Omoomi-Vahid Dastjerdi paper proves all gcd(n,k)\ensuremath{\geq}3 and selected gcd-two cases, and discusses earlier Zhu-Shao even/gcd-one classes. It does not establish all odd n\ensuremath{\geq}7 for k=2. The explicit odd k=2 construction is a candidate new infinite region of GPAll; its proof replaces earlier incomplete k=2 and k=3 arguments.

\begin{samepage}
Known covers of Petersen and K4, selected even odd-k classes and earlier periodic families provide prior constructions. Historical novelty of the remaining parameter regions and reduction capabilities requires an exact theorem-by-theorem comparison; introducing a definition does not by itself establish a new obstruction-removal theorem.

\end{samepage}
An optional strong-five-colour hypothesis may be false in light of the Deng-Liu-Feng 2025 infinite cubic family requiring six colours. Its exact graph assumptions require full-text comparison. An implication from that hypothesis can remain logically valid, even if the hypothesis is false; it would then not be a route proving the conjecture. This concern does not invalidate the explicit GP and flower endpoints.

\subsection{Proof scope and open premises}
The formal statements verify the finite constructions, infinite-family constructions and conditional reduction implications as typed. The five infinite structural premises are not proved globally, and the general DMS conjecture remains open. A selected endpoint audit does not imply that every theorem in the source portfolio has the same axiom classification.

The supplementary mathematical exposition gives the generalized-Petersen periods, seams and exceptions for 1\ensuremath{\leq}k\ensuremath{\leq}15, with 519 finite instances and 20,612 representative windows checked separately. It also explains the I/II reduction, two-/three-cut gluing, prescribed matching-bit induction and universal host/marking census bridge. The base12, simple14 and b14d finite assemblies replace the earlier partial census data. The five infinite structural premises remain open.

The explicit construction scope excludes all k\ensuremath{\geq}16 and does not settle unrestricted cubic or subcubic star six-colourability.

\subsection{Formal verification}
A qualified body-identical import-only composite route checks all 228 scientific source bodies through 221 identity-checked prior outputs and seven new source checks. Four import headers differ from the frozen originals; statements, proofs, options and comments are unchanged. The original-source run remains partial at 119/228. Its original organizer audit produced 120/123 outputs and failed because Star6Equiv was not imported; a separate supplement adds that import and checks all 123 unchanged requests. Among these selected declarations, 120 use standard kernel axioms and three are axiom-free; none uses admitted or unknown axioms. MGraph.k4subdiv\_star6 in the imported StarCert finite certificate uses native\_decide outside the selected set, so the entire portfolio is not asserted standard-only. The five open structural premises and the general conjecture remain open. Exact source/dependency identities and the separate audit records are supplied in the companion repository.

\begin{samepage}
\subsection{gp\_star5}
\begin{quote}\small\ttfamily theorem gp\_star5 (k n : Nat) (hk : 1 \ensuremath{\leq} k) (hK : k \ensuremath{\leq} 15) (hn : 2 * k + 1 \ensuremath{\leq} n) (hex : \ensuremath{\neg} (n = 3 \ensuremath{\wedge} k = 1)) :\newline     \ensuremath{\exists} c : Fin (gp n k).m \ensuremath{\to} Fin 5, (gp n k).Star 5 c\end{quote}
{\small Source: entries/dms-star6/artifact/lean/pack4/src/GPAll.lean:25.\par}

{\small Source SHA-256: e831e584837dc1e0a2e4ef26bf7875bebd3a18ad54deff326b3c68d1af970893\par}

\end{samepage}
\begin{samepage}
\subsection{gp\_family}
\begin{quote}\small\ttfamily theorem gp\_family (k n : Nat) (hk : 1 \ensuremath{\leq} k) (hK : k \ensuremath{\leq} 15) (hn : 2 * k + 1 \ensuremath{\leq} n) (hex : \ensuremath{\neg} (n = 3 \ensuremath{\wedge} k = 1)) :\newline     StarFam (gp n k) 5\end{quote}
{\small Source: entries/dms-star6/artifact/lean/pack4/src/GPAll.lean:45.\par}

{\small Source SHA-256: e831e584837dc1e0a2e4ef26bf7875bebd3a18ad54deff326b3c68d1af970893\par}

{\small\textbf{Sources and attribution:} \href{https://arxiv.org/html/2410.15024v1}{1. arXiv source: 2410.15024v1}; \href{https://arxiv.org/abs/1011.3376}{2. arXiv source: 1011.3376}; \href{https://doi.org/10.7151/dmgt.2195}{3. DOI: 10.7151/dmgt.2195}; \href{https://link.springer.com/article/10.1007/s40840-025-01875-9}{4. link.springer.com / s40840-025-01875-9}; \href{https://github.com/alejandrozu/ulam-opdp-difficulty-atlas}{5. alejandrozu/ulam-opdp-difficulty-atlas}; \href{https://github.com/alejandrozu/openmath-2026-judging/blob/d0ed487f487fa7ec814ff705ab55249983fe8707/OpenMath-Judging/review/entries/E08-DMS.txt}{6. alejandrozu/openmath-2026-judging / E08-DMS.txt}; \href{https://github.com/alejandrozu/openmath-2026-judging}{Proof source repository}.\par}

\end{samepage}
{\small \textbf{AI assistance.} Codex assisted literature checking, editorial preparation and analysis of the supplied formal artifacts. The human authors are responsible for checking the final mathematical claims, references and attribution.\par}

\end{document}
